\documentclass[11pt]{article}

% ----------- Font Packages ------------
\usepackage[T1]{fontenc}
\usepackage[utf8]{inputenc}

% ----------- Math Formatting ----------
\usepackage{amsmath, amssymb, amsthm}
\usepackage{dsfont}
\usepackage[most]{tcolorbox}
\tcbuselibrary{theorems}

% Cross-referencing and hyperlinks
\usepackage{hyperref}
\usepackage{cleveref}

% ----------- Layout and Spacing -------
\usepackage[letterpaper, top=0.75in, bottom=0.75in, left=1in, right=1in, heightrounded]{geometry}
\usepackage{setspace}
\onehalfspacing

% ----------- Misc Enhancements --------
\usepackage{microtype}
\usepackage{enumitem}
\usepackage{fancyhdr}
\pagestyle{fancy}
\fancyhead[L]{Homework Three}

\usepackage{bookmark}

\hypersetup{
    colorlinks=true,
    linkcolor=blue,
    urlcolor=blue,
    citecolor=blue
}

\title{Algebra Homework: Exponents and Order of Operations}
\author{Name: \underline{\hspace{3cm}}}
\date{Due: Nov 14}

\begin{document}
\maketitle

% ---------------------------------------------------
\section*{Foundational Questions}

\begin{enumerate}
    \item Explain the difference between \(3^2\) and \(2^3\).  
    Why do switching the exponent and the base change the result?

    \item What does it mean for an exponent to apply only to the number or expression it is directly attached to?  
    Give an example where parentheses change the value of an exponent expression.

    \item If you see an expression like  
    \[
        4 + 2^3 \times (5 - 1),
    \]  
    which part do you evaluate first? Explain your reasoning using PEMDAS.
\end{enumerate}

\vspace{1cm}

% ---------------------------------------------------
\section*{Parentheses and Exponents Practice}

Work step-by-step. Simplify the parentheses first, then apply exponents, then multiply or divide.

\[
\textbf{(1)} \quad (6 - 2)^2 + 3
\]

\vspace{1.4cm}

\[
\textbf{(2)} \quad 4 \times (3^2 - 7)
\]

\vspace{1.4cm}

\[
\textbf{(3)} \quad 10 - \big( 2^3 + (4 - 1) \big)
\]

\vspace{1.8cm}

% ---------------------------------------------------
\newpage
\section*{Mixed Exponent Practice}

Simplify each expression:

\begin{enumerate}[label=\textbf{(\arabic*)}]
    \item \( 2^4 + 3^2 \)
    \item \( (5 - 1)^2 \times 2 \)
    \item \( 9 + (2^3 - 1) \)
    \item \( 4 \times (1 + 3^2) \)
    \item \( (2 + 3)^2 - (4 - 1) \)
\end{enumerate}

\vspace{1cm}

% ---------------------------------------------------
\section*{Word Problem Challenge}

A character in Roblox builds a tall tower.  
The base height of the tower is \((8 - 3)^2\) blocks.  
Another player destroys \( \frac{1}{4} \) of the tower.  
Then the character rebuilds \(2^3\) blocks on top.

\begin{enumerate}[label=\textbf{(\alph*)}]
    \item Write an expression that represents the final height of the tower.
    \item Simplify the expression fully.
\end{enumerate}

\end{document}
